\subsection{Mixed-Chain Availability, Recovery, and CI}
The fresh 1.9.0 Linux ARM64 campaign executes both the three-package
npm $\to$ Docker $\to$ Docker chain and the four-package running example.
Build outputs carry fixture content through OCI archives; private Verdaccio
and Docker registries distinguish local output from publication~\cite{oci,buildkit}.
Each topology repeats five configurations six times in seeded order: mixed and
serialized availability, a missing barrier, refused publication, and failed building.

\begin{table}[t]
\centering
\caption{Fresh 1.9.0 executions. Designed failures are retained controls;
checks are verifier predicates, not independent samples.}
\label{tab:mixed}
\small
\setlength{\tabcolsep}{2pt}
\begin{tabular}{@{}lrrl@{}}
\toprule
Scope & Runs & Checks & Initial outcomes \\
\midrule
Three-package chain & 30 & 522 & 12 pass; 18 fail \\
Four-package chain & 30 & 798 & 12 pass; 18 fail \\
CI entry point & 3 & 81 & 3 pass \\
\bottomrule
\end{tabular}
\end{table}

Per topology, 12 repairable failures completed with unchanged versions; 24
completed cases then reran without build or publication. Six missing-barrier
controls were not repaired. Enabled barriers placed app building after base
publication; without one, the app failed against the unavailable image. In all
six mixed three-package runs, base building overlapped npm publication.

In the four-package mixed runs, UI building overlapped delayed assets
publication. All six base builds finished before UI publication started. This
is a valid local-output schedule: base building need not await UI publication,
but overlap is not required. Intervals use a shared monotonic clock~\cite{pytime}.
An initial campaign attempt incorrectly required
strict overlap and failed its oracle. We corrected that predicate before a
full rerun from fresh state; a focused oracle test accepts completion before publication
starts and rejects a build starting after publication ends. The correction and
earlier failure are retained, not counted as implementation evidence.

The three CI repetitions kept the workflow entry point and shared configuration
unchanged while adding packages. They selected required tests through the graph,
blocked release stages for an injected or missing test, propagated a specification
change to its consumer, and observed a final no-op release rerun. These fixtures
support selected ordering, recovery, and CI-integration paths on one Linux ARM64
environment, not a universal schedule, reliability rate, throughput result, or
independent adoption study.
